
\subsection{滤过 $A_{n,K}$ 同态的构造}
\label{subsec:7-2-9}
\liftIndex{I256}%

本小节证明命题 4.6.14 与定理 4.6.1
的 $A_{n,K}$ 版本（暂不计符号）。
先固定记号。

回忆 $\psi\colon M\to M'$ 是辛微分同胚，
且 $\psi(L)=L'$。
取正数列 $\mathcal K_\lambda,\delta_\lambda$，
满足当 $\lambda\to\infty$ 时
$\mathcal K_\lambda\nearrow\infty$、
$\delta_\lambda\searrow0$。
对每个 $\lambda$（相应地，对 $\lambda'$），
将命题~\ref{prop:7-2-35}
用于 $\mathcal K=\mathcal K_\lambda$、
$\delta=\delta_\lambda$
（相应地，用于
$\mathcal K_{\lambda'},\delta_{\lambda'}$），
得到可数集
\[
\mathcal X_{1,g}(L,\lambda)
\quad(g\le\mathcal K_\lambda),
\qquad
\mathcal X_{2,g'}(L',\lambda')
\quad(g'\le\mathcal K_{\lambda'}),
\]
及各自的多值截面
\[
\mathfrak s_{L,J,k,\beta,\vec P},
\qquad
\mathfrak s_{2,L',J',k,\beta',\vec P'}.
\]
它们分别定义在
\[
\mathcal M^{\mathrm{main}}_{k+1}
(L,J;\beta;\vec P),
\qquad
\mathcal M^{\mathrm{main},\mathfrak d'}_{k+1}
(L',J';\beta';\vec P')
\]
上；当
$\|(\mathfrak d,\beta)\|\le\mathcal K_\lambda$
（相应地
$\|(\mathfrak d',\beta')\|\le\mathcal K_{\lambda'}$）
时，其零集按命题~\ref{prop:7-2-35}
三角剖分。
下文只考虑 $\lambda\le\lambda'$ 的情形。
这里沿用第 4.6 节的记号，
$J$ 和 $J'$ 分别用于 $L$ 与 $L'$ 的模空间。
第二个多值截面的下标 $2$ 表示它关联于
$\mathcal X_{2,g'}(L',\lambda')$；
后面的命题 7.2.100 将把没有这个下标的
$\mathfrak s_{L',J',\mathfrak d',k,\beta',\vec P'}$
留给包含该集合的另一可数集
$\mathcal X_{3,g'}(L';(\lambda,\lambda'))$。

令
\[
\begin{aligned}
\mathcal X_{1,(\mathcal K_\lambda)}(L,\lambda)
&=\bigcup_{g\le\mathcal K_\lambda}
\mathcal X_{1,g}(L,\lambda),\\
\mathcal X_{2,(\mathcal K_{\lambda'})}(L',\lambda')
&=\bigcup_{g'\le\mathcal K_{\lambda'}}
\mathcal X_{2,g'}(L',\lambda').
\end{aligned}
\]
它们分别生成
\[
\begin{aligned}
C_{1,+}(L,\lambda)
&:=C_{1,(\mathcal K_\lambda)}(L,\lambda),\\
C_{2,+}(L',\lambda')
&:=C_{2,(\mathcal K_{\lambda'})}(L',\lambda').
\end{aligned}
\]
依命题~\ref{prop:7-2-35}，
可定义 $\mathfrak m^{\lambda,\mathrm{geo}}_{k,\beta}$
和 $\mathfrak m^{\lambda',\mathrm{geo}}_{k,\beta'}$；
只要 $A_\infty$ 关系的两端都有定义，
它们就满足该关系。
又依 (7.2.91)、(7.2.92)，
选取 $g_0(\lambda),g_1(\lambda)$
及 $g_0(\lambda'),g_1(\lambda')$，并令
\[
\begin{aligned}
C_1(L,\lambda)
&=C_{1,(g_1(\lambda))}(L,\lambda),&
C_{1,0}(L,\lambda)
&=C_{1,(g_0(\lambda))}(L,\lambda),\\
C_2(L',\lambda')
&=C_{2,(g_1(\lambda'))}(L',\lambda'),&
C_{2,0}(L',\lambda')
&=C_{2,(g_0(\lambda'))}(L',\lambda').
\end{aligned}
\]
由命题~\ref{prop:7-2-78}，
$C_1(L,\lambda)$ 与 $C_2(L',\lambda')$
分别具有滤过
$A_{n(\lambda),K(\lambda)}$
与 $A_{n(\lambda'),K(\lambda')}$ 结构
$\mathfrak m^\lambda_k$
与 $\mathfrak m^{\lambda'}_k$，且
\[
\lim_{\lambda\to\infty}n(\lambda)
=\lim_{\lambda\to\infty}K(\lambda)=\infty.
\]

构造滤过 $A_\infty$ 同态的策略
与代数结构的构造相似。
先定义满足下列性质 7.2.95 的可数集
$\mathcal X_{3,g'}(L';(\lambda,\lambda'))$，并置
\[
\mathcal X_{3,(\mathcal K'_{\lambda,\lambda'})}
(L';(\lambda,\lambda'))
=\bigcup_{g'\le\mathcal K'_{\lambda,\lambda'}}
\mathcal X_{3,g'}(L';(\lambda,\lambda')).
\]
它生成
\[
C_{3,+}(L';(\lambda,\lambda'))
:=C_{3,(\mathcal K'_{\lambda,\lambda'})}
(L';(\lambda,\lambda')),
\]
其中选取
$\mathcal K'_{\lambda,\lambda'}
\le\mathcal K_\lambda,\mathcal K_{\lambda'}$。
随后在有定义的范围内构造
``$A_{n'(\lambda),K'(\lambda)}$ 同态''
\[
\mathfrak f^{(\lambda,\lambda'),\mathrm{geo}}\colon
C_{1,+}(L;\lambda)
\longrightarrow C_{3,+}(L';(\lambda,\lambda')),
\]
并在目标上构造运算
$\mathfrak m^{(\lambda,\lambda'),\mathrm{geo}}$。
只要关系两端都有定义，
后者满足 $A_{n'(\lambda),K'(\lambda)}$ 关系，
且 $\mathfrak f^{(\lambda,\lambda'),\mathrm{geo}}$
同 $\mathfrak m^{\lambda,\mathrm{geo}}$、
$\mathfrak m^{(\lambda,\lambda'),\mathrm{geo}}$
满足相应的同态关系。

再仿照第~\ref{subsec:7-2-7} 小节，
在
$C_3(L';(\lambda,\lambda'))
\subset C_{3,+}(L';(\lambda,\lambda'))$
上取得 $A_{n'(\lambda),K'(\lambda)}$ 结构
$\mathfrak m^{(\lambda,\lambda')}$；
命题 7.2.126 将给出同阶同态
\[
\mathfrak f^{(\lambda,\lambda')}\colon
(C_1(L,\lambda),\mathfrak m^\lambda)
\longrightarrow
(C_3(L';(\lambda,\lambda')),
\mathfrak m^{(\lambda,\lambda')}).
\]
取 $\lambda_0<\lambda_1<\cdots\to\infty$，
并使
$(n'(\lambda_i),K'(\lambda_i))\to(\infty,\infty)$。
依命题 7.2.132 取极限，
将得到 $C_1(L,\lambda_0)$、
$C_3(L';(\lambda_0,\lambda_0'))$ 上的
滤过 $A_\infty$ 结构，
以及所需的滤过 $A_\infty$ 同态
\[
\mathfrak f\colon C_1(L,\lambda_0)
\longrightarrow C_3(L';(\lambda_0,\lambda_0')).
\]

现在落实上述方案。
先给出光滑奇异单形集
$\mathcal X_{3,g'}(L';(\lambda,\lambda'))$
应满足的条件。
对 $L$ 上的奇异单形 $P$，
记 $\psi(P)$ 为其在辛微分同胚
$\psi\colon M\to M'$ 下的像。

\begin{statement}{性质}{7.2.95}\label{prop:7-2-95}
\noindent\textup{(7.2.96.0)}\quad
$\mathcal X_{3,00}(L';(\lambda,\lambda'))$
是
$\mathcal X_{2,(\mathcal K_{\lambda'})}(L',\lambda')
\cup\psi(\mathcal X_{1,(\mathcal K_\lambda)}(L,\lambda))$
的并。

\noindent\textup{(7.2.96.1)}\quad
$\mathcal X_{3,g'}(L';(\lambda,\lambda'))$
是 $L'$ 上光滑奇异单形的可数集。

\noindent\textup{(7.2.96.2)}\quad
包含映射
$C_{3,(g')}(L';(\lambda,\lambda');\mathbb Q)
\to S(L';\mathbb Q)$
诱导上同调同构；
这里左端是由
$\mathcal X_{3,(g')}(L';(\lambda,\lambda'))$
在 $\mathbb Q$ 上生成的复形。

\noindent\textup{(7.2.96.3)}\quad
集合
$\mathcal X_{3,g'}(L';(\lambda,\lambda'))$
中每个元素的所有面，
均属于
\[
\mathcal X_{3,(g')}(L';(\lambda,\lambda'))
\]
这一可数集。
\end{statement}

条件 (7.2.96.1)--(7.2.96.3)
是滤过 $A_{n,K}$ 同态情形下
性质~\ref{prop:7-2-27} 的对应条件。
需注意，$\psi(P)$（其中
$P\in\mathcal X_{1,(\mathcal K_\lambda)}(L,\lambda)$）
可能与某个
$P'\in\mathcal X_{2,(\mathcal K_{\lambda'})}(L',\lambda')$
相同，因而 (7.2.96.0) 的并未必不交。

若 $P'\in\mathcal X_{3,g'}(L';(\lambda,\lambda'))$，
称它在这一构造中的\emph{代}为 $g'$。
由 (7.2.96.0)，
$\mathcal X_{2,(\mathcal K_{\lambda'})}(L',\lambda')$
中的单形在这里代为零；
它可能不同于构造 $C_2(L';\Lambda_{0,\mathrm{nov}})$
上的滤过 $A_{n,K}$ 结构时采用的代。
当前的同态构造忽略后一个代，
并以新指定的代定义装饰模空间
$\mathcal M^{\mathrm{main},\mathfrak d'}_{k+1}
(L',J';\beta;\vec P')$。
另一方面，$L$ 上固定的可数生成链空间
$\mathcal X_{1,(\mathcal K_\lambda)}(L,\lambda)$
在同态构造中不再变动，
所以不在 $L$ 的模空间记号中显写装饰；
相应多值截面记作
$\mathfrak s_{L,J,k,\beta,\vec P}$，
而不是
$\mathfrak s_{L,J,k,\beta,\vec P}$。

再说明构造同态所用的模空间。
定义 4.6.10 的时间有序积模空间，
此处简记为
\[
\mathcal M^{\mathrm{main}}_{k+1}
(\{J_\rho\}_\rho:\beta;\operatorname{top}(\rho);\vec P)
:=\mathcal M^{\mathrm{main}}_{k+1}
(M',L',\{J_\rho\}_\rho:\beta;\operatorname{top}(\rho);\vec P).
\]
其中 $\vec P=(P_1,\ldots,P_k)$
均取自固定的
$\mathcal X_{1,(\mathcal K_\lambda)}(L,\lambda)$。
其评价映射为
\begin{equation}
\operatorname{ev}_0\colon
\mathcal M^{\mathrm{main}}_{k+1}
(\{J_\rho\}_\rho:\beta;
\operatorname{top}(\rho);\vec P)
\longrightarrow L'.
\tag{7.2.97}\label{eq:7-2-97}
\end{equation}

命题 7.2.100 将构造仓西映射的扰动
（即多值截面）
$\mathfrak s^{\operatorname{top}(\rho)}_{\beta,k,\vec P}$，
并定义链
\[
\mathfrak f^{\mathrm{geo}}_{k,\beta}(\vec P)
=\operatorname{ev}_{0*}\left(
\mathcal M^{\mathrm{main}}_{k+1}
(\{J_\rho\}_\rho:\beta;
\operatorname{top}(\rho);\vec P)
^{\mathfrak s^{\operatorname{top}(\rho)}_{\beta,k,\vec P}}
\right).
\]
如命题 4.6.14 之后所述，
该模空间边界的余维一分量
由 (4.6.16.1)--(4.6.16.5) 描述；
其中 (4.6.16.1) 与 (4.6.16.2)
的贡献相消。

先看 (4.6.16.3)，它可写成
\begin{equation}
\mathcal M^{\mathrm{main}}_{k+\ell-m+1}
\bigl(\{J_\rho\}_\rho:\beta_2;
\operatorname{top}(\rho);
(P_1,\ldots,P_{\ell-1},Q^0_{\ell,m},
P_{m+1},\ldots,P_k)\bigr),
\tag{7.2.98.1}\label{eq:7-2-98-1}
\end{equation}
其中 $\beta=\beta_1+\beta_2$，且
\begin{equation}
Q^0_{\ell,m}
=\operatorname{ev}_{0*}\bigl(
\mathcal M^{\mathrm{main}}_{m-\ell+2}
(L,J;\beta_1;(P_\ell,\ldots,P_m))\bigr).
\tag{7.2.98.2}\label{eq:7-2-98-2}
\end{equation}
式~\eqref{eq:7-2-98-2} 中的模空间
用于定义奇异链
$\mathfrak m^{J,\mathfrak s}_{m-\ell+1}
(P_\ell,\ldots,P_m)$；
这里显写 $J,\mathfrak s$ 以表明依赖。
定义这一运算时，
该模空间的扰动已经固定。
注意 $P_i$ 始终取自不再变动的
$\mathcal X_{1,(\mathcal K_\lambda)}(L,\lambda)$。

再看 (4.6.16.4)，它按实际的 $m$ 个输入
可写成下式。%
\translatornote{原书 (7.2.99) 把末端写为
$R_{\ell_{m-1},\ell_m}$，并取
$\ell_m=k+1$、$\sum_{i=1}^{m-1}\beta_{1,i}$；
但外层 $\mathcal M_{m+1}$ 有 $m$ 个输入，
后文 (7.2.109) 也逐项使用
$\ell_{m+1}$ 与 $\sum_{i=1}^{m}\beta_{1,i}$。
这里统一补足第 $m$ 项；参见疑误 A-0443。}
\begin{equation}
\mathcal M^{\mathrm{main}}_{m+1}
(L',J';\beta_2;
R^0_{\ell_1,\ell_2},\ldots,
R^0_{\ell_m,\ell_{m+1}}),
\tag{7.2.99.1}\label{eq:7-2-99-1}
\end{equation}
其中
$\beta=\beta_2+\sum_{i=1}^{m}\beta_{1,i}$、
$1=\ell_1\le\ell_2\le\cdots\le\ell_{m+1}=k+1$，且
\begin{equation}
R^0_{\ell_i,\ell_{i+1}}
=\operatorname{ev}_{0*}\left(
\mathcal M^{\mathrm{main}}_{\ell_{i+1}-\ell_i+1}
(\{J_\rho\}_\rho:\beta_{1,i};
\operatorname{top}(\rho);
(P_{\ell_i},\ldots,P_{\ell_{i+1}-1}))
\right).
\tag{7.2.99.2}\label{eq:7-2-99-2}
\end{equation}
式~\eqref{eq:7-2-99-1} 的模空间
用于定义 $C_3(L';\Lambda_{0,\mathrm{nov}})$
上的滤过 $A_{n,K}$ 结构，
而式~\eqref{eq:7-2-99-2} 的模空间
用于定义相应的滤过 $A_{n,K}$ 同态。
不同于 (7.2.98)，
这两个模空间的扰动都尚未固定。

下面给出证明命题 4.6.14
所需的主要技术结果。

\begin{proposition}{7.2.100}\label{prop:7-2-100}
取正数列
$\mathcal K_\lambda,\delta_\lambda$，
使当 $\lambda\to\infty$ 时
$\mathcal K_\lambda\nearrow\infty$、
$\delta_\lambda\searrow0$。
由本小节开头的构造，
已有 $\mathcal X_{1,g}(L,\lambda)$、
$\mathcal X_{2,g'}(L',\lambda')$
以及多值截面
$\mathfrak s_{L,J,k,\beta,\vec P}$、
$\mathfrak s_{2,L',J',\mathfrak d',k,\beta,\vec P'}$
等。假设 $\lambda\le\lambda'$。
沿用上述记号，另取正数列
$\delta'_{\lambda,\lambda'}$、
$\mathcal K'_{\lambda,\lambda'}$，使
\[
\begin{gathered}
\mathcal K'_{\lambda,\lambda'}\le
\mathcal K_\lambda,\mathcal K_{\lambda'},
\qquad
\delta'_{\lambda,\lambda'}\ge
\delta_\lambda,\delta_{\lambda'},\\
\mathcal K'_{\lambda,\lambda'}\nearrow\infty,
\qquad
\delta'_{\lambda,\lambda'}\searrow0
\quad(\lambda\to\infty).
\end{gathered}
\]
只要 $\lambda,\lambda'$ 充分大，
就存在下列对象：
当 $\|(\mathfrak d',\beta)\|
\le\mathcal K'_{\lambda,\lambda'}$ 时，
在 $L'$ 的装饰模空间上的多值截面
$\mathfrak s_{L',J',\mathfrak d',k,\beta,\vec P'}$；
当 $\|\beta\|+k\le
\mathcal K'_{\lambda,\lambda'}$ 时，
分段光滑多值截面
$\mathfrak s^{\operatorname{top}(\rho)}_{\beta,k,\vec P}$；
以及 $g'\le\mathcal K'_{\lambda,\lambda'}$
时的可数集
$\mathcal X_{3,g'}(L';(\lambda,\lambda'))$。
它们满足下列性质。

\noindent\textup{(7.2.101.1)}\quad
$\mathcal X_{3,g'}(L';(\lambda,\lambda'))$
满足性质~\ref{prop:7-2-95}。

\noindent\textup{(7.2.101.2)}\quad
当 $\vec P'=(P_1',\ldots,P_k')$，
$P_i'\in\mathcal X_{3,g_i'}(L';(\lambda,\lambda'))$
时，
$\mathfrak s_{L',J',\mathfrak d',k,\beta,\vec P'}$
是
$\mathcal M^{\mathrm{main},\mathfrak d'}_{k+1}
(L',J';\beta;\vec P')$
上与零截面横截的多值截面。
若各 $P_i'$ 还属于
$\mathcal X_{2,g_i'}(L',\lambda')$
（$g_i'\le\mathcal K_{\lambda'}$），
则它与先前为 $L'$ 上的滤过结构构造的
$\mathfrak s_{2,L',J',\mathfrak d',k,\beta,\vec P'}$
一致。
此外，对取自
$\mathcal X_{1,g}(L,\lambda)$
（$g\le\mathcal K_\lambda$）的
$\vec P=(P_1,\ldots,P_k)$，
$\mathfrak s^{\operatorname{top}(\rho)}_{\beta,k,\vec P}$
是
$\mathcal M^{\mathrm{main}}_{k+1}
(\{J_\rho\}_\rho:\beta;
\operatorname{top}(\rho);\vec P)$
上与零截面横截的多值截面。

\noindent\textup{(7.2.101.3)}\quad
若 $\|(\mathfrak d',\beta)\|=g'$，
则链
\[
\operatorname{ev}_{0*}\left(
\mathcal M^{\mathrm{main},\mathfrak d'}_{k+1}
(L',J';\beta;\vec P')
^{\mathfrak s_{L',J',\mathfrak d',k,\beta,\vec P'}}
\right)
\]
可分解为
$\mathcal X_{3,(g')}(L';(\lambda,\lambda'))$
中的元素。

\noindent\textup{(7.2.101.4)}\quad
若 $\|\beta\|+k=g'$，则链
\[
\operatorname{ev}_{0*}\left(
\mathcal M^{\mathrm{main}}_{k+1}
(\{J_\rho\}_\rho:\beta;
\operatorname{top}(\rho);\vec P)
^{\mathfrak s^{\operatorname{top}(\rho)}_{\beta,k,\vec P}}
\right)
\]
可分解为同一个
$\mathcal X_{3,(g')}(L';(\lambda,\lambda'))$
中的元素。

\noindent\textup{(7.2.101.5)}\quad
$\mathfrak s_{L,J,k,\beta,\vec P}$
与 $\mathfrak s^{\operatorname{top}(\rho)}_{\beta,k,\vec P}$
按下文相容条件 7.2.102 相容；
$\mathfrak s_{L',J',\mathfrak d',k,\beta,\vec P'}$
与 $\mathfrak s^{\operatorname{top}(\rho)}_{\beta,k,\vec P}$
按相容条件 7.2.104 相容。

\noindent\textup{(7.2.101.6)}\quad
$\mathfrak s_{L,J,k,\beta,\vec P}$ 与
$\mathfrak s^{\operatorname{top}(\rho)}_{\beta,k,\vec P}$
的零集，均位于对应原仓西映射零集的
$\delta'_{\lambda,\lambda'}$-邻域内；
$\mathfrak s_{L',J',\mathfrak d',k,\beta,\vec P'}$
与
$\mathfrak s^{\operatorname{top}(\rho)}_{\beta,k,\vec P}$
的零集也满足同一邻域条件。
\end{proposition}

为何 $\mathfrak s^{\operatorname{top}(\rho)}_{\beta,k,\vec P}$
只能要求分段光滑而非光滑，
见注 7.2.122。
在 (7.2.101.2) 中有
\[
\mathcal X_{2,g'}(L',\lambda')
\subset\mathcal X_{3,00}(L';(\lambda,\lambda'))
\subset\mathcal X_{3,0}(L';(\lambda,\lambda'))
\subset\mathcal X_{3,g'}(L';(\lambda,\lambda')).
\]
命题~\ref{prop:7-2-100}
对每对充分大的 $(\lambda,\lambda')$ 均成立。
为证明它，以下暂固定这样的参数，
到证明末尾为止省去记号中的 $\lambda,\lambda'$，并置
\[
\begin{aligned}
\mathcal X_1(L)
&=\mathcal X_{1,(\mathcal K_\lambda)}(L,\lambda),\\
\mathcal X_2(L')
&=\mathcal X_{2,(\mathcal K_{\lambda'})}(L',\lambda'),\\
\mathcal X_3(L')
&=\mathcal X_{3,(\mathcal K'_{\lambda,\lambda'})}
(L';(\lambda,\lambda')).
\end{aligned}
\]
先说明 (7.2.101.5) 中的相容性。
$L'$ 上诸
$\mathfrak s_{L',J',\mathfrak d',k,\beta,\vec P'}$
之间的相容性与 (7.2.36.4) 完全相同，
不再重复。

现在说明
$\mathfrak s^{\operatorname{top}(\rho)}_{\beta,k,\vec P}$
在边界上与其他多值截面的相容性。
其模空间的边界由 (7.2.98)、(7.2.99)，
或某个输入变为 $\partial P_i$ 的模空间
\[
\mathcal M^{\mathrm{main}}_{k+1}
(\{J_\rho\}_\rho:\beta;
\operatorname{top}(\rho);
(P_1,\ldots,\partial P_i,\ldots,P_k))
\]
描述。
最后一种情形的相容性
与相容条件~\ref{cond:7-2-38} 相同。

先看 (7.2.98)。
$L$ 上
$\mathcal M^{\mathrm{main}}_{m-\ell+2}
(L,J;\beta_1;(P_\ell,\ldots,P_m))$
的扰动，在定义 $\mathfrak m^{\mathfrak s,J}$
时已经固定。
因此有链
\[
Q_{\ell,m}
=\operatorname{ev}_{0*}\left(
\mathcal M^{\mathrm{main}}_{m-\ell+2}
(L,J;\beta_1;(P_\ell,\ldots,P_m))
^{\mathfrak s_{L,J,m-\ell+1,\beta_1,
(P_\ell,\ldots,P_m)}}
\right),
\]
它属于 $\mathcal X_1(L)$。
命题~\ref{prop:7-2-100}
又为以 $Q_{\ell,m}$ 作一个输入的
外层模空间
\[
\mathcal M^{\mathrm{main}}_{k+\ell-m+1}
(\{J_\rho\}_\rho:\beta_2;
\operatorname{top}(\rho);
(P_1,\ldots,P_{\ell-1},Q_{\ell,m},
P_{m+1},\ldots,P_k))
\]
提供多值截面。

\begin{statement}{相容条件}{7.2.102}
\label{cond:7-2-102}
\liftIndex{I241}
下列两个多值截面的直和
\[
\begin{aligned}
&\mathfrak s^{\operatorname{top}(\rho)}_{\beta_2,
k-m+\ell,(P_1,\ldots,P_{\ell-1},Q_{\ell,m},
P_{m+1},\ldots,P_k)}\\
&\qquad\oplus
\mathfrak s_{L,J,m-\ell+1,\beta_1,
(P_\ell,\ldots,P_m)}
\end{aligned}
\]
等于
$\mathfrak s^{\operatorname{top}(\rho)}_{\beta,k,\vec P}$
在相应边界上的拉回。见图~\ref{fig:7-2-6}。
\end{statement}

\par\medskip\noindent\begin{minipage}{\textwidth}
\centering
\setlength{\unitlength}{0.72mm}
\begin{picture}(128,64)
\linethickness{0.5pt}
\put(45,34){\circle{36}}
\put(81,34){\circle{36}}
\put(45,34){\makebox(0,0){$\beta_1$}}
\put(81,34){\makebox(0,0){$\beta_2$}}
\put(63,34){\circle*{1.6}}
\put(99,34){\vector(1,0){12}}
\put(113,33){$\scriptstyle z_0$}
\put(30,54){$\scriptstyle P_\ell$}
\put(54,54){$\scriptstyle\cdots$}
\put(33,11){$\scriptstyle P_m$}
\put(76,55){$\scriptstyle P_1,\ldots$}
\put(87,9){$\scriptstyle P_{m+1},\ldots,P_k$}
\put(54,4){\vector(-1,0){24}}
\put(17,1){$\scriptstyle Q_{\ell,m}$}
\qbezier[20](30,54)(12,34)(30,14)
\qbezier[20](60,54)(73,34)(60,14)
\qbezier[20](69,54)(57,34)(69,14)
\qbezier[20](93,54)(113,34)(93,14)
\end{picture}
\par{\renewcommand{\thefigure}{7.2.6}\refstepcounter{figure}
\normalfont\bfseries 图~\thefigure\label{fig:7-2-6}}
\end{minipage}\par\medskip

此处的拉回与相容条件~\ref{cond:7-2-44}
中的拉回相同。
严格地说，
相容条件 7.2.102 也应如
相容条件~\ref{cond:7-2-44} 那样
给出更细致的陈述；
两者的细节处理相同，
所以这里只写简式。
以后出现类似情形时不再重复这一提醒。

下面再说明与 (7.2.99) 对应的相容条件。
命题~\ref{prop:7-2-100}
为每个分支提供
$\mathfrak s^{\operatorname{top}(\rho)}_{\beta_1,
m-\ell+1,(P_\ell,\ldots,P_m)}$，
其定义域按输入数应为下式。%
\translatornote{原书 (7.2.103.1) 印作
$\mathcal M_{m-\ell}^{\mathrm{main}}$
且把类写为 $\beta_2$；
但这里有 $m-\ell+1$ 个输入，
紧前的多值截面也属 $\beta_1$。
按下一式 (7.2.103.2) 的一般型别，
应为 $\mathcal M_{m-\ell+2}^{\mathrm{main}}$
与 $\beta_1$；参见疑误 A-0444。}
\begin{equation}
\mathcal M^{\mathrm{main}}_{m-\ell+2}
(\{J_\rho\}_\rho:\beta_1;
\operatorname{top}(\rho);
(P_\ell,\ldots,P_m)).
\tag{7.2.103.1}\label{eq:7-2-103-1}
\end{equation}
这里须用后面证明的引理 7.2.107
保证 $\|\beta_1\|+m-\ell+1\le\mathcal K$。
定义
\begin{equation}
\begin{aligned}
R_{\ell_i,\ell_{i+1}}
=\operatorname{ev}_{0*}\bigl(&
\mathcal M^{\mathrm{main}}_{\ell_{i+1}-\ell_i+1}
(\{J_\rho\}_\rho:\beta_{1,i};\\[-2pt]
&\operatorname{top}(\rho);
(P_{\ell_i},\ldots,P_{\ell_{i+1}-1}))
^{\mathfrak s^{\operatorname{top}(\rho)}_{\beta_{1,i},
\ell_{i+1}-\ell_i,(P_{\ell_i},\ldots,
P_{\ell_{i+1}-1})}}
\bigr).
\end{aligned}
\tag{7.2.103.2}\label{eq:7-2-103-2}
\end{equation}
置
\[
\mathfrak d'^{\max}(i)
=\|\beta_{1,i}\|+\ell_{i+1}-\ell_i.
\]
由 (7.2.101.4)，链 $R_{\ell_i,\ell_{i+1}}$
可分解为
$\mathcal X_{3,(\mathfrak d'^{\max}(i))}(L')$
中的元素。
为免记号繁重，
以下把 $R_{\ell_i,\ell_{i+1}}$
视为其中一个元素，
其代记为 $g_i\le\mathfrak d'^{\max}(i)$。
置
\[
\begin{aligned}
\vec R&=(R_{\ell_1,\ell_2},\ldots,
R_{\ell_m,\ell_{m+1}}),\\
\mathfrak d'&\colon\{1,\ldots,m\}
\longrightarrow\{0,1,\ldots,\mathcal K'_{\lambda,\lambda'}\},
\qquad\mathfrak d'(i)=g_i.
\end{aligned}
\]
命题~\ref{prop:7-2-100} 又提供
$\mathcal M^{\mathrm{main},\mathfrak d'}_{m+1}
(L',J';\beta_2;\vec R)$
上的多值截面
$\mathfrak s_{L',J',\mathfrak d',m,\beta_2,\vec R}$；
这里仍需使用下文的引理~\ref{lem:7-2-107}。
其中 $m$ 为输入数，
与疑误 A-0443 的裁决一致。

\begin{statement}{相容条件}{7.2.104}
\label{cond:7-2-104}
\liftIndex{I241}
下列多值截面的直和
\[
\begin{aligned}
&\mathfrak s_{L',J',\mathfrak d',m,\beta_2,\vec R}
\oplus\bigoplus_{i=1}^{m}
\mathfrak s^{\operatorname{top}(\rho)}_{\beta_{1,i},
\ell_{i+1}-\ell_i,
(P_{\ell_i},\ldots,P_{\ell_{i+1}-1})}
\end{aligned}
\]
等于
$\mathfrak s^{\operatorname{top}(\rho)}_{\beta,k,\vec P}$
在相应边界上的拉回。见图~\ref{fig:7-2-7}。
\end{statement}

\par\medskip\noindent\begin{minipage}{\textwidth}
\centering
\setlength{\unitlength}{0.72mm}
\begin{picture}(128,108)
\linethickness{0.5pt}
\put(88,52){\circle{40}}
\put(75,84){\circle{32}}
\put(52,52){\circle{32}}
\put(75,20){\circle{32}}
\put(88,52){\makebox(0,0){$\beta_2$}}
\put(75,84){\makebox(0,0){$\scriptstyle\beta_{1,1}$}}
\put(52,52){\makebox(0,0){$\scriptstyle\beta_{1,2}$}}
\put(75,20){\makebox(0,0){$\scriptstyle\beta_{1,m}$}}
\put(108,52){\vector(1,0){11}}
\put(120,51){$\scriptstyle z_0$}
\put(16,84){\circle{12}}
\put(16,52){\circle{12}}
\put(16,20){\circle{12}}
\put(7,94){$\scriptstyle R_{\ell_1,\ell_2}$}
\put(7,62){$\scriptstyle R_{\ell_2,\ell_3}$}
\put(5,30){$\scriptstyle R_{\ell_m,\ell_{m+1}}$}
\put(22,84){\vector(1,0){37}}
\put(22,52){\vector(1,0){14}}
\put(22,20){\vector(1,0){37}}
\put(62,101){$\scriptstyle P_{\ell_1},\ldots,P_{\ell_2-1}$}
\put(34,70){$\scriptstyle P_{\ell_2},\ldots$}
\put(62,1){$\scriptstyle P_{\ell_m},\ldots,P_{\ell_{m+1}-1}$}
\put(77,68){\circle*{1.4}}
\put(68,52){\circle*{1.4}}
\put(77,36){\circle*{1.4}}
\qbezier[18](60,97)(45,85)(60,71)
\qbezier[18](37,68)(23,52)(37,36)
\qbezier[18](60,33)(45,20)(60,7)
\end{picture}
\par{\renewcommand{\thefigure}{7.2.7}\refstepcounter{figure}
\normalfont\bfseries 图~\thefigure\label{fig:7-2-7}}
\end{minipage}\par\medskip

这样便精确陈述了命题~\ref{prop:7-2-100}。
其证明是命题~\ref{prop:7-2-35} 证明的直接对应物：
归纳构造 $L'$ 上与时间有序积模空间上的多值截面
及 $\mathcal X_{3,g}(L')$，
使之满足相容条件~\ref{cond:7-2-102}
和~\ref{cond:7-2-104}。
附录 A1 的定理 A1.23
（即 \textbf{[FuOn99II]} 的引理 3.11）
可把归纳假设在边界上给出的多值截面延拓。
与引理~\ref{lem:7-2-55} 的证明相同，
还可证明相应的引理。
因此只需说明定义多值截面所用的次序，
并核对归纳能否继续。

构造多值截面的次序如下。

\begin{definition}{7.2.105}\label{def:7-2-105}
\noindent\textup{(7.2.106.1)}\quad
若下列情形之一成立，则
$\mathfrak s_{L',J',\mathfrak d'_1,k'_1,\beta_1,\vec P'_{(1)}}
<\mathfrak s_{L',J',\mathfrak d'_2,k'_2,\beta_2,\vec P'_{(2)}}$：

\noindent\textup{(7.2.106.1.1)}\quad
$\|(\mathfrak d'_1,\beta_1)\|<\|(\mathfrak d'_2,\beta_2)\|$；

\noindent\textup{(7.2.106.1.2)}\quad
上述两范数相等，且
$k'_1+\|\beta_1\|<k'_2+\|\beta_2\|$；

\noindent\textup{(7.2.106.1.3)}\quad
上述两范数及 $k'_i+\|\beta_i\|$ 分别相等，且
$\|\beta_1\|<\|\beta_2\|$；

\noindent\textup{(7.2.106.1.4)}\quad
上述两范数、$k'_i+\|\beta_i\|$ 及
$\|\beta_i\|$ 分别相等，且
$\sum_i\deg P'_{(1),i}>
\sum_i\deg P'_{(2),i}$。

\noindent\textup{(7.2.106.2)}\quad
若下列情形之一成立，则
$\mathfrak s^{\operatorname{top}(\rho)}_{\beta_1,k_1,\vec P_{(1)}}
<\mathfrak s^{\operatorname{top}(\rho)}_{\beta_2,k_2,\vec P_{(2)}}$：

\noindent\textup{(7.2.106.2.1)}\quad
$k_1+\|\beta_1\|<k_2+\|\beta_2\|$；

\noindent\textup{(7.2.106.2.2)}\quad
上述两值相等，且 $\|\beta_1\|<\|\beta_2\|$；

\noindent\textup{(7.2.106.2.3)}\quad
上述两值及 $\|\beta_i\|$ 分别相等，且
$\sum_i\deg P_{(1),i}>
\sum_i\deg P_{(2),i}$。

\noindent\textup{(7.2.106.3)}\quad
若下列情形之一成立，则
$\mathfrak s_{L',J',\mathfrak d'_1,k'_1,\beta_1,\vec P'_{(1)}}
<\mathfrak s^{\operatorname{top}(\rho)}_{\beta_2,k_2,\vec P_{(2)}}$：

\noindent\textup{(7.2.106.3.1)}\quad
$\|(\mathfrak d'_1,\beta_1)\|<k_2+\|\beta_2\|$；

\noindent\textup{(7.2.106.3.2)}\quad
上述两值相等，且
$k'_1+\|\beta_1\|<k_2+\|\beta_2\|$；

\noindent\textup{(7.2.106.3.3)}\quad
上述两值及 $k'_1+\|\beta_1\|$ 分别相等，且
$\|\beta_1\|<\|\beta_2\|$。

\noindent\textup{(7.2.106.4)}\quad
若下列情形之一成立，则
$\mathfrak s^{\operatorname{top}(\rho)}_{\beta_1,k_1,\vec P_{(1)}}
<\mathfrak s_{L',J',\mathfrak d'_2,k'_2,\beta_2,\vec P'_{(2)}}$：

\noindent\textup{(7.2.106.4.1)}\quad
$k_1+\|\beta_1\|<\|(\mathfrak d'_2,\beta_2)\|$；

\noindent\textup{(7.2.106.4.2)}\quad
上述两值相等，且
$k_1+\|\beta_1\|<k'_2+\|\beta_2\|$。
（事实上，此情形不会发生。）

\noindent\textup{(7.2.106.4.3)}\quad
上述两值及 $k'_2+\|\beta_2\|$ 分别相等，且
$\|\beta_1\|<\|\beta_2\|$。
\end{definition}

下面证明引理 7.2.107 和 7.2.120，
它们表明命题~\ref{prop:7-2-100}
的证明可以依照上述次序 $<$ 归纳进行。

\begin{lemma}{7.2.107}\label{lem:7-2-107}
设 $R_{\ell_i,\ell_{i+1}}$、
$\mathfrak d'^{\max}$、$\beta$、$\beta_2$、
$\beta_{1,i}$ 等如 (7.2.103) 所定义，则
\begin{align}
\mathfrak s_{L',J',\mathfrak d',m,\beta_2,\vec R}
&<\mathfrak s^{\operatorname{top}(\rho)}_{\beta,k,\vec P},
\tag{7.2.108.1}\label{eq:7-2-108-1}\\
\mathfrak s^{\operatorname{top}(\rho)}_{\beta_{1,i},
\ell_{i+1}-\ell_i,
(P_{\ell_i},\ldots,P_{\ell_{i+1}-1})}
&<\mathfrak s^{\operatorname{top}(\rho)}_{\beta,k,\vec P}
\quad\text{对每个 $i$ 成立}.
\tag{7.2.108.2}\label{eq:7-2-108-2}
\end{align}
\end{lemma}

\begin{proof}
先证 (7.2.108.1)。不妨设
$\mathfrak d'=\mathfrak d'^{\max}$，因为
\[
\mathfrak s_{L',J',\mathfrak d',m,\beta_2,\vec R}
\le\mathfrak s_{L',J',\mathfrak d'^{\max},m,\beta_2,\vec R}.
\]
我们有
\begin{align}
\mathfrak d'(j)&=\|\beta_{1,j}\|+\ell_{j+1}-\ell_j,
\tag{7.2.109.1}\label{eq:7-2-109-1}\\
\sum_{j=1}^{m}(\ell_{j+1}-\ell_j)&=k,
\tag{7.2.109.2}\label{eq:7-2-109-2}\\
\beta_2+\sum_{j=1}^{m}\beta_{1,j}&=\beta.
\tag{7.2.109.3}\label{eq:7-2-109-3}
\end{align}

\par\medskip\noindent\begin{minipage}{\textwidth}
\centering
\setlength{\unitlength}{0.72mm}
\begin{picture}(120,94)
\linethickness{0.5pt}
\put(66,42){\circle{40}}
\put(66,42){\makebox(0,0){$\beta_2$}}
\put(87,42){\vector(1,0){12}}
\put(101,41){$\scriptstyle z_0$}
\put(53,74){\circle{30}}
\put(53,74){\makebox(0,0){$\beta_{1,1}$}}
\put(32,42){\circle{28}}
\put(32,42){\makebox(0,0){$\beta_{1,2}$}}
\put(78,10){\circle{20}}
\put(78,10){\makebox(0,0){$\beta_{1,m}$}}
\put(43,91){$\scriptstyle\ell_2-\ell_1$}
\put(61,59){\circle*{1.4}}
\put(46,42){\circle*{1.4}}
\put(73,22){\circle*{1.4}}
\qbezier[15](38,86)(27,73)(39,62)
\qbezier[15](20,55)(9,43)(20,29)
\qbezier[15](35,25)(49,8)(64,5)
\end{picture}
\par{\renewcommand{\thefigure}{7.2.8}\refstepcounter{figure}
\normalfont\bfseries 图~\thefigure\label{fig:7-2-8}}
\end{minipage}\par\medskip

取 $j_0$ 使
$\mathfrak d'(j_0)=\max\mathfrak d'$。
于是
\begin{equation}
\begin{aligned}
\|(\mathfrak d',\beta_2)\|
&=\|\beta_2\|+\mathfrak d'(j_0)+m\\
&=\|\beta_2\|+\|\beta_{1,j_0}\|
+(\ell_{j_0+1}-\ell_{j_0})+m.
\end{aligned}
\tag{7.2.110}\label{eq:7-2-110}
\end{equation}
由 (7.2.49.2) 与 (7.2.109.3) 得
\begin{equation}
\|\beta\|\ge\|\beta_2\|
+\sum_{j=1}^{m}\|\beta_{1,j}\|+m.
\tag{7.2.111}\label{eq:7-2-111}
\end{equation}
再由 (7.2.110)、(7.2.111) 得
\begin{equation}
\|(\mathfrak d',\beta_2)\|
\le\|\beta\|-\sum_{j\ne j_0}\|\beta_{1,j}\|
+(\ell_{j_0+1}-\ell_{j_0}).
\tag{7.2.112}\label{eq:7-2-112}
\end{equation}
由稳定性，
$\|\beta_{1,j}\|+\ell_{j+1}-\ell_j\ge0$；
等号仅当
$\|\beta_{1,j}\|=0$、
$\ell_{j+1}=\ell_j$ 且
$\beta_{1,j}\ne\beta_0$ 时成立。
故由 (7.2.109.2) 与 (7.2.112)，或者
\begin{equation}
\|(\mathfrak d',\beta_2)\|<\|\beta\|+k,
\tag{7.2.113.1}\label{eq:7-2-113-1}
\end{equation}
或者
\begin{equation}
\|(\mathfrak d',\beta_2)\|=\|\beta\|+k,
\qquad
\|\beta_{1,j}\|=0,\quad\ell_{j+1}=\ell_j
\quad(j\ne j_0).
\tag{7.2.113.2}\label{eq:7-2-113-2}
\end{equation}
在后一情形下，(7.2.111) 也取等号。
若 (7.2.113.1) 成立，便符合 (7.2.106.3.1)，
从而结论成立。

现在设 (7.2.113.2) 成立。
由 (7.2.109.2) 与 (7.2.112)，有
\begin{equation}
\|(\mathfrak d',\beta_2)\|
=\|\beta\|+(\ell_{j_0+1}-\ell_{j_0})
=\|\beta\|+k.
\tag{7.2.114}\label{eq:7-2-114}
\end{equation}
又有
$m+\|\beta_2\|
=\|(\mathfrak d',\beta_2)\|-\mathfrak d'(j_0)$。
因此或者
\begin{equation}
\|(\mathfrak d',\beta_2)\|=\|\beta\|+k,
\quad m+\|\beta_2\|<k+\|\beta\|,
\tag{7.2.115.1}\label{eq:7-2-115-1}
\end{equation}
或者
\begin{equation}
\mathfrak d'(j_0)=0.
\tag{7.2.115.2}\label{eq:7-2-115-2}
\end{equation}
若 (7.2.115.1) 成立，便符合 (7.2.106.3.2)。
若 (7.2.115.2) 成立，则由 (7.2.109.1) 有
\begin{equation}
\|\beta_{1,j_0}\|=0,
\qquad \ell_{j_0+1}-\ell_{j_0}=0.
\tag{7.2.116}\label{eq:7-2-116}
\end{equation}
式 (7.2.116)、(7.2.113.2)、(7.2.109.2)
推出 $k=0$。
于是 (7.2.114) 给出
$\|(\mathfrak d',\beta_2)\|=\|\beta\|$。
由于 (7.2.111) 取等号，
(7.2.116) 与 (7.2.113.2) 还给出
$\|\beta\|=\|\beta_2\|+m$。
注意 $m>0$，故符合 (7.2.106.3.3)。
这就证毕 (7.2.108.1)。

再证 (7.2.108.2)。
由 (7.2.109.2)，以及
$\|\beta_{1,i}\|\le\|\beta\|$、
$\ell_{j+1}-\ell_j\ge0$，得
\begin{equation}
\ell_{i+1}-\ell_i+\|\beta_{1,i}\|
\le\|\beta\|+k.
\tag{7.2.117}\label{eq:7-2-117}
\end{equation}
若为严格不等式，则符合 (7.2.106.2.1)。
若为等式，则
\begin{equation}
\ell_{j+1}-\ell_j=0\quad(j\ne i),
\tag{7.2.118}\label{eq:7-2-118}
\end{equation}
且
\begin{equation}
\beta_2=\beta_0,
\qquad\beta_{1,j}=\beta_0\quad(j\ne i).
\tag{7.2.119}\label{eq:7-2-119}
\end{equation}
由稳定性，(7.2.118)、(7.2.119)
推出 $m=1$，即不存在 $j\ne i$。
但稳定性及 $m=1$ 又要求
$\beta_2\ne\beta_0$，
与 (7.2.119) 矛盾。
故 (7.2.108.2) 得证。
\end{proof}

\begin{lemma}{7.2.120}\label{lem:7-2-120}
在相容条件~\ref{cond:7-2-102} 的情形下，有
\[
\mathfrak s^{\operatorname{top}(\rho)}_{
\beta_2,k-m+\ell,
(P_1,\ldots,P_{\ell-1},Q_{\ell,m},P_{m+1},\ldots,P_k)}
<\mathfrak s^{\operatorname{top}(\rho)}_{\beta,k,\vec P}.
\]
\end{lemma}
\begin{proof}
$\beta_1+\beta_2=\beta$。
若 $\beta_1=\beta_0$，则 $\ell<m$，因而
$\|\beta_2\|+k-m+\ell<k+\|\beta\|$。
若 $\beta_1\ne\beta_0$，则
$\|\beta_2\|+1\le\|\beta\|$，
$k-m+\ell\le k+1$。
因此，或者
$\|\beta_2\|+k-m+\ell<\|\beta\|+k$，
或者两边相等且
$\|\beta_2\|<\|\beta\|$。
这分别对应 (7.2.106.2.1) 和 (7.2.106.2.2)。
\end{proof}

为研究
$\partial\mathcal M^{\mathrm{main},\mathfrak d'}_{k+1}
(L',J';\beta;\vec P')$，
还需一个同类引理，
但它恰好就是引理~\ref{lem:7-2-47}。

现在完成命题~\ref{prop:7-2-100} 的证明。
先定义第零代元素的集合
$\mathcal X_{3,0}(L')$。
取可数集
\[
\mathcal X_{3,00}(L')
=\mathcal X_2(L')\cup\psi(\mathcal X_1(L)).
\]
若时间有序积模空间
\[
\mathcal M^{\mathrm{main}}_{k+1}
(\{J_\rho\}_\rho:\beta;
\operatorname{top}(\rho);\vec P)
\]
满足 $k+\|\beta\|=0$，
则由引理~\ref{lem:7-2-45}，
或者 (i) $k=0$、$\beta\ne\beta_0$、
$\|\beta\|=0$，
或者 (ii) $k=1$、$\beta=\beta_0$。
稳定性排除了 (ii)。
在 (i) 中，模空间
$\mathcal M^{\mathrm{main}}_{0+1}
(\{J_\rho\}_\rho:\beta;
\operatorname{top}(\rho);\mathbf 1)$
没有边界。
固定其上与零截面横截的多值截面
$\mathfrak s^{\operatorname{top}(\rho)}_{\beta,0,\mathbf 1}$，
再将
\[
\operatorname{ev}_{0*}\left(
\mathcal M^{\mathrm{main}}_{0+1}
(\{J_\rho\}_\rho:\beta;
\operatorname{top}(\rho);\mathbf 1)
^{\mathfrak s^{\operatorname{top}(\rho)}_{\beta,0,\mathbf 1}}
\right)
\]
作单纯剖分，把它视为光滑奇异链；
其剖分中尚未属于
$\mathcal X_{3,00}(L')$ 的奇异单形，
也加入其中。这样便得到
$\mathcal X_{3,0}(L')$。

还须指出，当 $\|(\mathfrak d',\beta)\|=0$，
若
$\mathcal M^{\mathrm{main},\mathfrak d'}_{k+1}
(L',J';\beta;\vec P')$ 非空，
则 $k=0$、$\beta\ne\beta_0$、
$\|\beta\|=0$。
因此
$\mathcal M^{\mathrm{main},\mathfrak d'}_{0+1}
(L',J';\beta;\mathbf 1)$ 的扰动
已在上一小节固定。

其余归纳步骤与命题~\ref{prop:7-2-35}
的证明相同。
这里还需一个类似引理~\ref{lem:7-2-55}
的引理，同样可依原引理的证明给出。
命题~\ref{prop:7-2-100} 得证。

现在解释为什么命题~\ref{prop:7-2-100}
只能得到分段光滑、而未必光滑的
$\mathfrak s^{\operatorname{top}(\rho)}_{\beta,k,\vec P}$。
构造多值截面时，先在带仓西结构的空间的
各边界分量上给定多值截面，
使它们在角上相容，再延拓到边界邻域，
然后应用定理 A1.23。
这一步实际使用下列引理。
回忆：定义在 $\mathbb R^n$ 子集上的函数称为光滑，
是指它为某个定义在该子集邻域上的光滑函数的限制。

\begin{lemma}{7.2.121}\label{lem:7-2-121}
令
\[
\begin{aligned}
V&=\{(x_1,\ldots,x_n)\in\mathbb R^n
\mid x_i\ge0,\ i=1,\ldots,n\},\\
V_i&=V\cap\{(x_1,\ldots,x_n)\in\mathbb R^n
\mid x_i=0\}.
\end{aligned}
\]
设 $f_i\colon V_i\to\mathbb R$ 为光滑函数，
且在 $V_i\cap V_j$ 上 $f_i=f_j$。
则存在光滑函数 $f\colon V\to\mathbb R$，
使 $f|_{V_i}=f_i$。
\end{lemma}
\begin{proof}
通过加上常数，不妨设
$f_i(0)=0$（$i=1,\ldots,n$）。
对 $x=(x_1,\ldots,x_n)\in V_i$，定义
\[
f'_i(x)=f_i(x)-
\sum_{j\ne i}f_i(0,\ldots,0,
\underset{j}{x_j},0,\ldots,0).
\]
这些 $f'_i$ 仍在 $V_i\cap V_j$ 上相等，
而且在 $j\ne i$ 时还满足
\[
f'_i(0,\ldots,0,\underset{j}{x_j},0,\ldots,0)=0.
\]
由于对 $k\ne i$，
\[
\sum_{j\ne i}
f_i(0,\ldots,0,\underset{j}{x_j},0,\ldots,0)
+f_k(0,\ldots,0,\underset{i}{x_i},0,\ldots,0)
\]
是 $V$ 上的光滑函数，
我们不妨假设原来的 $f_i$ 在
$j\ne i$ 时就满足
\[
f_i(0,\ldots,0,\underset{j}{x_j},0,\ldots,0)=0.
\]
再对不同指标 $j,k$ 的两两组合重复这一过程，
可进一步化到
\[
f_i(0,\ldots,0,\underset{j}{x_j},0,\ldots,0,
\underset{k}{x_k},0,\ldots,0)=0
\qquad(j,k\ne i).
\]
如此归纳，最终化到对每个 $i$ 都有 $f_i=0$，
此时结论显然成立。
\end{proof}

\begin{remark}{7.2.122}\label{rem:7-2-122}
另一方面，引理~\ref{lem:7-2-121}
不能直接推广到命题~\ref{prop:7-2-100}
证明中出现的多边形。
设 $P\subset\mathbb R^n$ 为凸多边形，
并把 $\partial P$ 分解为 $\bigcup_aP_a$。
若 $P_a$ 上的光滑函数 $f_a$
在 $P_a\cap P_b$ 上彼此相等，
一般仍未必存在 $P$ 上的光滑函数
$f\colon P\to\mathbb R$
使每个 $f|_{P_a}=f_a$。

回到命题~\ref{prop:7-2-100}，
考察
$\mathcal M^{\mathrm{main}}_{k+1}
(\{J_\rho\}_\rho:\beta;
\operatorname{top}(\rho);\vec P)$
上的多值截面
$\mathfrak s^{\operatorname{top}(\rho)}_{\beta,k,\vec P}$。
该模空间的边界分解为模空间的纤维积，
记分解中的分量为 $\{\mathcal M_\alpha\}_\alpha$。
这一分解由 (7.1.52) 诱导。
如注~\ref{rem:7-1-69} 所述，
(7.1.52) 并不是该模空间
（的仓西邻域）按不同维数的边界和角
所作的标准分解。
所以 $\{\mathcal M_\alpha\}_\alpha$
也不是边界和角的标准分解。
即使各 $\mathcal M_\alpha$ 上给定的多值截面
在交集处相容，也不能直接用
引理~\ref{lem:7-2-121} 延拓它们。
这就是没有断言
$\mathfrak s^{\operatorname{top}(\rho)}_{\beta,k,\vec P}$
光滑的原因。

不过仍可构造分段光滑的多值截面。
将上述模空间的仓西邻域分解为
$\{\mathcal N_\gamma\}_\gamma$，
使各 $\mathcal N_\gamma$ 都是带角流形，
而每个 $\mathcal M_\alpha$
都是某个 $\mathcal N_\gamma$
作为带角流形的标准分解中的分量；
例如可以作单纯剖分。
利用引理~\ref{lem:7-2-121}，
可把边界上给定的、在每个
$\mathcal M_\alpha$ 上均光滑的多值截面
延拓到边界邻域，使它在每个
$\mathcal N_\gamma$ 上光滑。
由此得到分段光滑多值截面。

就定义虚基本链而言，
分段光滑多值截面可与光滑多值截面同样使用：
把与零截面横截的多值截面之零集三角剖分，
即可由之得到虚基本链。
横截的分段光滑多值截面之零集确有三角剖分。
\end{remark}

由命题~\ref{prop:7-2-100}，定义
\[
\mathcal X_{3,(\mathcal K'_{\lambda,\lambda'})}
(L';(\lambda,\lambda'))
=\bigcup_{g'\le\mathcal K'_{\lambda,\lambda'}}
\mathcal X_{3,g'}(L';(\lambda,\lambda')),
\]
它生成
\[
C_{3,+}(L';(\lambda,\lambda'))
:=C_{3,(\mathcal K'_{\lambda,\lambda'})}
(L';(\lambda,\lambda')).
\]
在这个空间上，定义运算
$\mathfrak m^{(\lambda,\lambda'),\mathrm{geo}}$：
\begin{equation}
\mathfrak m^{(\lambda,\lambda'),\mathrm{geo}}_{k,\beta}
(\vec P')=
\operatorname{ev}_{0*}\left(
\mathcal M^{\mathrm{main},\mathfrak d'}_{k+1}
(L',J';\beta;\vec P')
^{\mathfrak s_{L',J',\mathfrak d',k,\beta,\vec P'}}
\right).
\tag{7.2.123}\label{eq:7-2-123}
\end{equation}
还定义一个部分有定义的滤过
``$A_{n'(\lambda),K'(\lambda)}$ 同态''
\[
\mathfrak f^{(\lambda,\lambda'),\mathrm{geo}}
\colon C_{1,+}(L,\lambda)
\longrightarrow C_{3,+}(L';(\lambda,\lambda')),
\]
其分量为
\begin{align}
\mathfrak f^{(\lambda,\lambda'),\mathrm{geo}}_{k,\beta}
(\vec P)
&=\operatorname{ev}_{0*}\left(
\mathcal M^{\mathrm{main}}_{k+1}
(\{J_\rho\}_\rho:\beta;
\operatorname{top}(\rho);\vec P)
^{\mathfrak s^{\operatorname{top}(\rho)}_{\beta,k,\vec P}}
\right),
\tag{7.2.124.1}\label{eq:7-2-124-1}\\
\mathfrak f^{(\lambda,\lambda'),\mathrm{geo}}_{1,\beta_0}
(P)&=\psi(P).
\tag{7.2.124.2}\label{eq:7-2-124-2}
\end{align}
这里 $n'(\lambda),K'(\lambda)$
由命题~\ref{prop:7-2-78}
对 $\mathcal K=\mathcal K'_{\lambda,\lambda'}$
给出，而后者取自命题~\ref{prop:7-2-100}。

只要 $A_{n'(\lambda),K'(\lambda)}$
关系两端都有定义，
$\mathfrak m^{(\lambda,\lambda'),\mathrm{geo}}$
就满足该关系。
只要同态关系两端都有定义，
$\mathfrak f^{(\lambda,\lambda'),\mathrm{geo}}$、
$\mathfrak m^{(\lambda,\lambda'),\mathrm{geo}}$、
$\mathfrak m^{\lambda,\mathrm{geo}}$
就满足滤过 $A_{n'(\lambda),K'(\lambda)}$
同态的定义式。

下面仿照第~\ref{subsec:7-2-7} 小节
修正这些运算。
先用命题~\ref{prop:7-2-78}
取得 $g_1(\lambda,\lambda')$、
$g_0(\lambda,\lambda')$，并置
\[
\begin{aligned}
C_3(L';(\lambda,\lambda'))
&=C_{3,(g_1(\lambda,\lambda'))}
(L';(\lambda,\lambda')),\\
C_{3,0}(L';(\lambda,\lambda'))
&=C_{3,(g_0(\lambda,\lambda'))}
(L';(\lambda,\lambda')).
\end{aligned}
\]
由此在 $C_3(L';(\lambda,\lambda'))$
上得到滤过 $A_{n'(\lambda),K'(\lambda)}$ 结构
$\mathfrak m^{(\lambda,\lambda')}_{k,\beta}$，满足
\begin{equation}
\mathfrak m^{(\lambda,\lambda')}_{k,\beta}
=\mathfrak m^{(\lambda,\lambda'),\mathrm{geo}}_{k,\beta}
\quad\text{在 }C_{3,0}(L';(\lambda,\lambda'))\text{ 上}.
\tag{7.2.125}\label{eq:7-2-125}
\end{equation}
回忆此结构由
$\mathfrak m^{(\lambda,\lambda'),\mathrm{geo}}_{k,\beta}$
及 $C_{3,+}(L';(\lambda,\lambda'))$
上的算子 $\mathfrak G_k$ 构造。
由构造有
\[
\mathfrak m^{(\lambda,\lambda'),\mathrm{geo}}_{k,\beta}
=\mathfrak m^{\lambda',\mathrm{geo}}_{k,\beta}
\quad\text{在 }
C_2(L',\lambda')\subset C_{3,+}(L';(\lambda,\lambda'))
\text{ 上}.
\]
参见 (7.2.101.2)。
因此可以选取 $\mathfrak G_k$，使
\[
\mathfrak m^{(\lambda,\lambda')}_{k,\beta}
=\mathfrak m^{\lambda'}_{k,\beta}
\quad\text{在 }C_2(L',\lambda')\text{ 上}.
\]

现在证明下列命题。
\begin{proposition}{7.2.126}\label{prop:7-2-126}
若 $\lambda'>C(\lambda)$，
则存在滤过 $A_{n'(\lambda),K'(\lambda)}$ 同态
\[
\mathfrak f^{(\lambda,\lambda')}
\colon(C_1(L,\lambda),\mathfrak m^\lambda)
\longrightarrow
(C_3(L';(\lambda,\lambda')),
\mathfrak m^{(\lambda,\lambda')})
\]
使得
$\mathfrak f^{(\lambda,\lambda'),\mathrm{geo}}$
将 $C_{1,0}(L,\lambda)$ 映入
$C_3(L';(\lambda,\lambda'))$，且
\begin{equation}
\mathfrak f^{(\lambda,\lambda')}
=\mathfrak f^{(\lambda,\lambda'),\mathrm{geo}}
\quad\text{在 }C_{1,0}(L,\lambda)\text{ 上}.
\tag{7.2.127}\label{eq:7-2-127}
\end{equation}
\end{proposition}

\begin{proof}
令
$\mathfrak f^{(\lambda,\lambda')}
=\mathfrak f^{(\lambda,\lambda'),\mathrm{geo}}
\circ\mathfrak i$，
其中 $\mathfrak i$ 如 (7.2.84)，
合成按滤过 $A_\infty$ 同态的情形定义。
适当增大 $c(\lambda)$，
可使 $\mathfrak f^{(\lambda,\lambda')}$
在 $C_1(L,\lambda)$ 上有定义，
且其像落在 $C_{3,0}(L';(\lambda,\lambda'))$ 中。
在后一空间上，
$\mathfrak m^{\mathrm{geo}}$ 与 $\mathfrak m$ 一致。
所以由命题 7.2.78bis 和 (7.2.125)，
容易证明 $\mathfrak f^{(\lambda,\lambda')}$
是滤过 $A_{n'(\lambda),K'(\lambda)}$ 同态。
(7.2.127) 则由命题~\ref{prop:7-2-78}
与定义得到。
\end{proof}

命题~\ref{prop:7-2-78} 与
命题~\ref{prop:7-2-126}
利用第 5.4 节的技巧，把
$\mathfrak m^{\mathrm{geo}}$、
$\mathfrak f^{\mathrm{geo}}$
修正为 $\mathfrak m$、$\mathfrak f$。
后面若干小节还会用到相同技巧；
由于方法一样，不再重复细节。

\begin{remark}{7.2.128}\label{rem:7-2-128}
若 $(M,L)=(M',L')$、$\psi=\operatorname{id}$、
$\lambda=\lambda'$、$\mathfrak s=\mathfrak s'$、
$J=J'$、
$\mathcal X_{1,s}(L,\lambda)
=\mathcal X_{2,s}(L',\lambda')$，
而定义运算 $\mathfrak m$ 的多值截面也相同，
则可在命题~\ref{prop:7-2-100} 中
选取多值截面，使
$C_{3,+}(L';(\lambda,\lambda'))
=C_2(L',\lambda')$，并且
\[
\mathfrak f^{(\lambda,\lambda'),\mathrm{geo}}
=\mathfrak f^{(\lambda,\lambda')}
=\operatorname{id}.
\]
上述等式的双模版本
用于第 6.1.4 小节末尾完成定理 5.3.14 的证明。
见注 7.2.172(2)。
\end{remark}

现在把命题~\ref{prop:7-2-126}
的滤过 $A_{n'(\lambda),K'(\lambda)}$
结构及同态延拓为滤过 $A_\infty$
结构及同态。为此需要如下通用引理；
沿用第~\ref{subsec:7-2-6} 小节的记号。
\liftIndex{I257}%
\begin{lemma}{7.2.129}\label{lem:7-2-129}
设 $(n,K)<(n',K')$，
$C_1,C_2,C'_1,C'_2$
均为滤过 $A_{n',K'}$ 代数。
设
$\mathfrak h\colon C_1\to C_2$、
$\mathfrak h'\colon C'_1\to C'_2$
为滤过 $A_{n',K'}$ 同伦等价；
$\mathfrak g_{(1)}\colon C_1\to C'_1$
为滤过 $A_{n,K}$ 同态；
$\mathfrak g_{(2)}\colon C_2\to C'_2$
为滤过 $A_{n',K'}$ 同态。
假设 $\mathfrak g_{(2)}\circ\mathfrak h$
与 $\mathfrak h'\circ\mathfrak g_{(1)}$
在 $A_{n,K}$ 意义下同伦。
则存在滤过 $A_{n',K'}$ 同态
$\mathfrak g^+_{(1)}\colon C_1\to C'_1$，
使它作为滤过 $A_{n,K}$ 同态
与 $\mathfrak g_{(1)}$ 一致，且
$\mathfrak g_{(2)}\circ\mathfrak h$
与 $\mathfrak h'\circ\mathfrak g^+_{(1)}$
在滤过 $A_{n',K'}$ 意义下同伦。
\end{lemma}
\begin{proof}
证明与定理~\ref{thm:7-2-72} 的证明相似，
均依据障碍理论。
只需考虑如下两种相邻阶情形：
\begin{align}
(n,K)&<(n',K')=(n+1,K-1),
\tag{7.2.130.1}\label{eq:7-2-130-1}\\
(n,K)&=(n,0)<(n',K')=(0,n+1).
\tag{7.2.130.2}\label{eq:7-2-130-2}
\end{align}
在每种情形中，都有延拓
$\mathfrak g_{(1)}$ 的障碍类
\[
[o_{K',\beta}(\mathfrak g_{(1)})]
\in H\bigl(\operatorname{Hom}
(B_{K'}\overline C_1[1],\overline C'_1[1])\bigr).
\]
参见 (7.2.76.2)；这里
$\beta\in\mathfrak G$、$\|\beta\|=n'$，
见定义~\ref{def:7-2-61}。
障碍类的自然性，以及
$\mathfrak g_{(2)}$ 的存在，
表明此障碍消失，
故可构造滤过 $A_{n',K'}$ 同态
$\mathfrak g^+_{(1)}$。

为把
$\mathfrak g_{(2)}\circ\mathfrak h$
与 $\mathfrak h'\circ\mathfrak g_{(1)}$
之间的滤过 $A_{n,K}$ 同伦
延拓为它与
$\mathfrak h'\circ\mathfrak g^+_{(1)}$
之间的滤过 $A_{n',K'}$ 同伦，
还会遇到类似第 4.5 节所研究的
\[
H\bigl(\operatorname{Hom}
(B_{K'}\overline C_1[1],\overline{\mathfrak C}'_2[1])\bigr)
\]
中的障碍类；其中
$\overline{\mathfrak C}'_2$ 是
$[0,1]\times\overline C'_2$ 的模型，
见第 4.2.1 小节。
延拓同伦的障碍理论表明，
适当改变 $\mathfrak g^+_{(1)}$ 的选取，
即可消除这些障碍。
论证见下文注~\ref{rem:7-2-131}，
亦可比较第~\ref{subsec:7-2-6} 小节
证明末尾的论证。
引理得证。
\end{proof}

\begin{remark}{7.2.131}\label{rem:7-2-131}
这里详述同伦延拓的障碍理论论证。
设 $(n,K)<(n',K')$ 如 (7.2.130)，
$\mathfrak g^{(0)},\mathfrak g^{(1)}
\colon C\to C'$ 为滤过 $A_{n',K'}$ 同态，
$\mathfrak F\colon C\to\mathfrak C'$
为二者之间的滤过 $A_{n,K}$ 同伦。
我们要把 $\mathfrak F$ 延拓为
滤过 $A_{n',K'}$ 同伦。
在引理~\ref{lem:7-2-129} 的证明中，取
\[
C=C_1,\qquad C'=C'_2,\qquad
\mathfrak C'=\mathfrak C'_2,\qquad
\mathfrak g^{(0)}=\mathfrak g_{(2)}\circ\mathfrak h,
\quad\mathfrak g^{(1)}=\mathfrak h'\circ\mathfrak g_{(1)}.
\]
令 $\|\beta\|=n'$。
因 $\overline{\mathfrak C}'$
是 $[0,1]\times\overline C'$ 的模型，
可取
$\mathfrak F^I_{K',\beta}\colon
B_{K'}(\overline C[1])\to\overline{\mathfrak C}'$
使
\[
\overline{\operatorname{Eval}}_{r=r_0}
\circ\mathfrak F^I_{K',\beta}
=\mathfrak g^{(r_0)}_{K',\beta},
\qquad r_0=0,1.
\]
以
\[
o_{K',\beta}(\mathfrak F)
\in\operatorname{Hom}
(B_{K'}(\overline C[1]),\overline{\mathfrak C}'[1])
\]
表示把 $\mathfrak F$ 延拓为
滤过 $A_{n',K'}$ 同态时的障碍循环。
置
\[
\mathfrak K=\ker\left(
\overline{\operatorname{Eval}}_{r=0}
\oplus\overline{\operatorname{Eval}}_{r=1}
\colon\overline{\mathfrak C}'
\longrightarrow
\overline C'_{r=0}\oplus\overline C'_{r=1}
\right).
\]
易见
$o_{K',\beta}(\mathfrak F)
+\delta_1(\mathfrak F^I_{K',\beta})
\in\operatorname{Hom}(B_{K'}(\overline C[1]),
\mathfrak K[1])$；
这里 $\delta_1$ 由 $\overline{\mathfrak m}_1$ 诱导。
由于 $\overline{\operatorname{Eval}}_{r=0}$
是同伦等价，正合列
\[
\begin{aligned}
0&\longrightarrow
\operatorname{Hom}(B_{K'}(\overline C[1]),\mathfrak K[1])
\longrightarrow
\operatorname{Hom}(B_{K'}(\overline C[1]),
\overline{\mathfrak C}'[1])\\
&\longrightarrow
\operatorname{Hom}(B_{K'}(\overline C[1]),
\overline C'_{r=0}[1]\oplus\overline C'_{r=1}[1])
\longrightarrow0
\end{aligned}
\]
表明连接同态
\[
\delta_1^\sharp\colon
H^{p-1}\bigl(\operatorname{Hom}
(B_{K'}(\overline C[1]),\overline C'_{r=1}[1])\bigr)
\longrightarrow
H^p\bigl(\operatorname{Hom}
(B_{K'}(\overline C[1]),\mathfrak K[1])\bigr)
\]
是同构。
因此存在次数为零的
$\operatorname{Cor1}_{K',\beta}
\in\operatorname{Hom}
(B_{K'}(\overline C[1]),\overline C'_{r=1}[1])$，
使
\[
\delta_1(\operatorname{Cor1})=0,
\qquad
\delta_1^\sharp([\operatorname{Cor1}])
=[o_{K',\beta}(\mathfrak F)
+\delta_1(\mathfrak F^I_{K',\beta})].
\]
于是存在
$\operatorname{Cor2}_{K',\beta}
\in\operatorname{Hom}
(B_{K'}(\overline C[1]),\overline{\mathfrak C}'[1])$，
满足
\[
\overline{\operatorname{Eval}}_{r=0}\circ
\operatorname{Cor2}_{K',\beta}=0,
\qquad
\overline{\operatorname{Eval}}_{r=1}\circ
\operatorname{Cor2}_{K',\beta}
=\operatorname{Cor1}_{K',\beta}.
\]
还存在
$\operatorname{Cor3}_{K',\beta}
\in\operatorname{Hom}
(B_{K'}(\overline C[1]),\mathfrak K[1])$，
使
\[
o_{K',\beta}(\mathfrak F)
+\delta_1(\mathfrak F^I_{K',\beta})
=\delta_1(\operatorname{Cor2}_{K',\beta})
+\delta_1(\operatorname{Cor3}_{K',\beta}).
\]
置
\[
\begin{aligned}
\mathfrak F_{K',\beta}
&=\mathfrak F^I_{K',\beta}
-\operatorname{Cor2}_{K',\beta}
-\operatorname{Cor3}_{K',\beta},\\
\mathfrak g'^{(1)}_{K',\beta}
&=\mathfrak g^{(1)}_{K',\beta}
-\operatorname{Cor1}_{K',\beta}.
\end{aligned}
\]
于是 $\mathfrak g'^{(1)}$
是滤过 $A_{n',K'}$ 同态，
而 $\mathfrak F$ 是
$\mathfrak g^{(0)}$ 与 $\mathfrak g'^{(1)}$
之间的滤过 $A_{n',K'}$ 同伦。
\end{remark}

现在可以完成定理 4.6.1 的证明
（暂不计符号）。
取 $\lambda=\lambda'$，
已有滤过 $A_{n'(\lambda),K'(\lambda)}$ 同态
\[
\mathfrak f^{(\lambda,\lambda)}
\colon C_1(L,\lambda)
\longrightarrow C_3(L';(\lambda,\lambda)).
\]
此时也可把 $C_1(L,\lambda)$
视为滤过 $A_{n'(\lambda),K'(\lambda)}$ 代数。
起初它是滤过 $A_{n(\lambda),K(\lambda)}$ 代数，
其中 $n(\lambda),K(\lambda)$
由命题~\ref{prop:7-2-78}、
命题~\ref{prop:7-2-35}
根据 $\mathcal K_\lambda$ 决定；
而 $n'(\lambda),K'(\lambda)$
由命题~\ref{prop:7-2-126}、
命题~\ref{prop:7-2-100}
根据 $\mathcal K'_{\lambda,\lambda}$ 决定。
由于 $\mathcal K'_{\lambda,\lambda}
\le\mathcal K_\lambda$，
有
$(n'(\lambda),K'(\lambda))
\le(n(\lambda),K(\lambda))$。

令
\[
\mathfrak g^\lambda\colon
C_1(L,\lambda)\longrightarrow C_2(L',\lambda)
\]
为 $\mathfrak f^{(\lambda,\lambda)}$
与包含映射
$C_2(L',\lambda)\to C_3(L';(\lambda,\lambda))$
的同伦逆所合成的滤过
$A_{n'(\lambda),K'(\lambda)}$ 同态。
只须把 $\mathfrak g^{\lambda_0}$
延拓为滤过 $A_\infty$ 同态。
回忆当 $\lambda\to\infty$ 时，
$n'(\lambda)\nearrow\infty$、
$K'(\lambda)\nearrow\infty$。

\begin{proposition}{7.2.132}\label{prop:7-2-132}
存在 $n''(\lambda),K''(\lambda)$，
当 $\lambda\to\infty$ 时
二者均趋于无穷，
且满足：若 $\lambda<\lambda'$，
则存在滤过 $A_{n''(\lambda),K''(\lambda)}$
同伦等价
\[
\begin{aligned}
\mathfrak h^{1,(\lambda,\lambda')}
&\colon C_1(L,\lambda)\to C_1(L,\lambda'),\\
\mathfrak h^{2,(\lambda,\lambda')}
&\colon C_2(L',\lambda)\to C_2(L',\lambda'),
\end{aligned}
\]
使
$\mathfrak g^{\lambda'}\circ
\mathfrak h^{1,(\lambda,\lambda')}$
与
$\mathfrak h^{2,(\lambda,\lambda')}\circ
\mathfrak g^\lambda$
在 $A_{n''(\lambda),K''(\lambda)}$
意义下同伦。
\end{proposition}

命题~\ref{prop:7-2-132}
由下一小节的结果推出，
其证明在第 7.2.11 小节给出。
设 $\lambda<\lambda'$。
定理~\ref{thm:7-2-72} 表明，
可将 $C_1(L,\lambda)$、
$C_2(L',\lambda)$ 上的滤过
$A_{n'(\lambda),K'(\lambda)}$ 结构
延拓为滤过
$A_{n'(\lambda'),K'(\lambda')}$ 结构，
并把
$\mathfrak h^{1,(\lambda,\lambda')}$、
$\mathfrak h^{2,(\lambda,\lambda')}$
延拓为同阶同态；
仍用相同符号表示延拓后的对象。
再由引理~\ref{lem:7-2-129}
和命题~\ref{prop:7-2-132}，
滤过 $A_{n'(\lambda),K'(\lambda)}$
同态 $\mathfrak g^\lambda$
可延拓为滤过
$A_{n'(\lambda'),K'(\lambda')}$ 同态，
仍记为 $\mathfrak g^\lambda$，
使
$\mathfrak g^{\lambda'}\circ
\mathfrak h^{1,(\lambda,\lambda')}$
与
$\mathfrak h^{2,(\lambda,\lambda')}\circ
\mathfrak g^\lambda$
在 $A_{n''(\lambda),K''(\lambda)}$
意义下同伦。

最后仿照第~\ref{subsec:7-2-8} 小节末尾，
取 $\lambda_0<\lambda_1<\cdots\to\infty$，
便在 $C_1(L,\lambda_0)$、
$C_2(L',\lambda_0)$ 上得到
滤过 $A_\infty$ 结构及滤过 $A_\infty$ 同态
\[
\mathfrak g\colon
C_1(L,\lambda_0)\to C_2(L',\lambda_0).
\]
再与同伦等价
$C_2(L',\lambda_0)
\to C_3(L';(\lambda_0,\lambda_0))$
合成，得到所需的滤过 $A_\infty$ 同态
\[
\mathfrak f\colon C_1(L,\lambda_0)
\longrightarrow C_3(L';(\lambda_0,\lambda_0)).
\]
定理 4.6.1 得证。
